\section{Separation of Variables}
\label{sec:separation-of-variables}

\begin{example}[Separation of variables for the heat equation]
\label{ex:heat-equation-separation-of-variables}
\dcref{ch4:4.1-ex-1}
\group{Separation of Variables}
For a bounded smooth $U\subset\R^n$, seek $u(x,t)=v(t)w(x)$ in
\[
u_t-\Delta u=0\text{ in }U\times(0,\infty),\quad u|_{\partial U}=0.
\]
Separation gives $v'/v=\Delta w/w=\mu$, hence $v=de^{\mu t}$ and
$-\Delta w=\lambda w$ with $\mu=-\lambda$. Thus eigenpairs yield
\[
u(x,t)=\sum_k d_ke^{-\lambda_kt}w_k(x),
\]
whenever $g=\sum_kd_kw_k$ (finite or convergent countable expansion).
\end{example}

\begin{definition}[Eigenvalue and eigenfunction]
\label{def:eigenvalue-eigenfunction}
\dcref{ch4:4.1-def-1}
\group{Separation of Variables}
For bounded open $U$, $\lambda$ is an eigenvalue of $-\Delta$ with zero
boundary data if some nonzero $w$ solves
\[
-\Delta w=\lambda w\text{ in }U,\qquad w=0\text{ on }\partial U.
\]
Such a $w$ is an eigenfunction.
\end{definition}

\begin{remark}[Nonlinear extensions require linearity]
\label{rem:separation-of-variables-linearity-remark}
\uses{ex:heat-equation-separation-of-variables}
\dcref{ch4:4.1-rem-1}
\group{Separation of Variables}

The single separated mode follows from the ansatz itself; superposition of
modes requires linearity of the PDE.
\end{remark}

\begin{example}[Separation of variables for the porous medium equation]
\label{ex:porous-medium-separation-of-variables}
\dcref{ch4:4.1-ex-2}
\group{Similarity Solutions}
For $u_t-\Delta(u^\gamma)=0$, $\gamma>1$, seek $u=v(t)w(x)$. Then
\[
\frac{v'}{v^\gamma}=\mu=\frac{\Delta(w^\gamma)}{w}.
\]
Taking $w=|x|^\alpha$ gives
\[
\alpha=\frac{2}{\gamma-1},\qquad
\mu=\alpha\gamma(\alpha\gamma+n-2)>0,
\]
and, for $\lambda>0$,
\[
u(x,t)=((1-\gamma)\mu t+\lambda)^{1/(1-\gamma)}|x|^\alpha.
\]
\end{example}

\begin{remark}[Finite-time blow-up]
\label{rem:porous-medium-blowup-remark}
\uses{ex:porous-medium-separation-of-variables}
\dcref{ch4:4.1-rem-2}
\group{Similarity Solutions}

The preceding solution blows up at
\[
t_*=(\gamma-1)^{-1}\mu^{-1}
\]
for $x\ne0$.
\end{remark}

\begin{example}[Additive separation of variables for the Hamilton--Jacobi equation]
\label{ex:hamilton-jacobi-additive-separation}
\uses{ex:hamilton-jacobi-complete-integral}
\dcref{ch4:4.1-ex-3}
\group{Separation of Variables}

For $u_t+H(Du)=0$, the additive ansatz $u=w(x)+v(t)$ gives
$H(Dw)=\mu=-v'(t)$. Hence every solution of $H(Dw)=\mu$ produces
\[
u(x,t)=w(x)-\mu t+b.
\]
In particular, $w(x)=a\cdot x$ yields $u=a\cdot x-H(a)t+b$.
\end{example}

\section{Similarity Solutions}
\label{sec:similarity-solutions}

\subsection{Plane and Traveling Waves, Solitons}
\label{sec:plane-traveling-waves-solitons}

\begin{definition}[Traveling wave and plane wave]
\label{def:traveling-wave-plane-wave}
\dcref{ch4:4.2.1-def-1}
\group{Traveling Waves}
A traveling wave has the form
\[
u(x,t)=v(x-\sigma t),
\]
and a plane wave in $\R^n$ has the form
\[
u(x,t)=v(y\cdot x-\sigma t),
\]
with wavefront normal $y$ and velocity $\sigma/|y|$.
\end{definition}

\begin{example}[Exponential plane wave solutions of linear PDE]
\label{ex:exponential-plane-wave-solutions}
\uses{def:traveling-wave-plane-wave}
\dcref{ch4:4.2.1-ex-1}
\group{Traveling Waves}

For $u=e^{i(y\cdot x+\omega t)}$, the dispersion relations are:
\[
\begin{array}{c|c|c}
\text{PDE}&\omega&\text{mode}\\ \hline
u_t-\Delta u=0&i|y|^2&e^{iy\cdot x-|y|^2t}\\
u_{tt}-\Delta u=0&\pm|y|&e^{i(y\cdot x\pm|y|t)}\\
u_t+u_{xxx}=0&y^3&e^{i(yx+y^3t)}\\
iu_t+\Delta u=0&-|y|^2&e^{i(y\cdot x-|y|^2t)}.
\end{array}
\]
The heat modes dissipate, while wave, Airy, and Schrodinger modes are
dispersive (the Airy velocity is $y^2$).
\end{example}

\begin{example}[Solitons for the Korteweg--de Vries equation]
\label{ex:kdv-soliton-traveling-wave}
\uses{def:traveling-wave-plane-wave}
\dcref{ch4:4.2.1-ex-2}
\group{Traveling Waves}

For KdV, $u_t+6uu_x+u_{xxx}=0$, a traveling profile satisfies the
third-order ODE
\[
-\sigma v'+6vv'+v'''=0.
\]
Integrating once gives
\[
v''-\sigma v+3v^2=a,\qquad
\frac{(v')^2}{2}=-v^3+\frac\sigma2v^2+av+b.
\]
Decay at both infinities forces $a=b=0$, and the resulting soliton is
\[
u(x,t)=\frac\sigma2\sech^2\!\left(\frac{\sqrt\sigma}{2}(x-\sigma t-c)\right),
\qquad \sigma>0,\ c\in\R.
\]
\end{example}

\begin{remark}[Complete integrability of the KdV equation]
\label{rem:kdv-complete-integrability}
\uses{ex:kdv-soliton-traveling-wave}
\dcref{ch4:4.2.1-rem-1}
\group{Traveling Waves}

KdV is completely integrable; see \cite{Drazin1983} for the associated
solution-generating machinery.
\end{remark}

\begin{theorem}[Existence of a traveling wave for the bistable reaction-diffusion equation]
\label{thm:bistable-traveling-wave-existence}
\uses{def:traveling-wave-plane-wave}
\dcref{ch4:4.2.1-thm-1}
\group{Traveling Waves}

Let $f$ be smooth and, for some $0<a<1$, satisfy
\[
f(0)=f(a)=f(1)=0,\quad f<0\text{ on }(0,a),\quad f>0\text{ on }(a,1),
\]
\[
f'(0)<0,\quad f'(1)<0,\quad \int_0^1f(z)\,dz>0.
\]
Then $u_t-u_{xx}=f(u)$ admits a traveling wave
$u(x,t)=v(x-\sigma t)$ with $v(-\infty)=0$ and $v(+\infty)=1$.
\end{theorem}

\begin{proof}
\sketch The profile solves $v''+\sigma v'+f(v)=0$. Writing $w=v'$ gives
$v'=w$, $w'=-\sigma w-f(v)$, with saddle equilibria $(0,0)$ and $(1,0)$.
The energy $E=w^2/2+\int_0^vf$ obeys $\dot E=-\sigma w^2$. For large negative
$\sigma$, unstable and stable trajectories cross a transverse line in the
opposite order; continuity in $\sigma$ yields a value where they coincide,
forming a heteroclinic orbit.
\end{proof}

\begin{remark}[Uniqueness of the bistable traveling wave velocity]
\label{rem:bistable-traveling-wave-velocity-uniqueness}
\uses{thm:bistable-traveling-wave-existence, ex:kdv-soliton-traveling-wave}
\dcref{ch4:4.2.1-rem-2}
\group{Traveling Waves}

A refined phase-plane argument gives uniqueness of the admissible velocity
$\sigma$ for fixed $f$, unlike the KdV soliton family.
\end{remark}

\subsection{Similarity Under Scaling}
\label{sec:similarity-under-scaling}

\begin{example}[A scaling invariant solution of the porous medium equation: Barenblatt's solution]
\label{ex:porous-medium-scaling-invariant-barenblatt}
\uses{def:fundamental-solution-heat-equation}
\dcref{ch4:4.2.2-ex-1}
\group{Similarity Solutions}

For $u_t-\Delta(u^\gamma)=0$, seek
\[
u(x,t)=t^{-\alpha}v(xt^{-\beta}),\qquad \alpha+1=\alpha\gamma+2\beta.
\]
For radial $v(y)=w(|y|)$, imposing $\alpha=n\beta$ gives
\[
(r^{n-1}(w^\gamma)')'+\beta(r^nw)'=0,
\quad
w(r)=\left(b-\frac{\gamma-1}{2\gamma}\beta r^2\right)_+^{1/(\gamma-1)},
\]
with
\[
\alpha=\frac n{n(\gamma-1)+2},\qquad
\beta=\frac1{n(\gamma-1)+2}.
\]
Thus
\[
u(x,t)=t^{-\alpha}\left(b-\frac{\gamma-1}{2\gamma}\beta\frac{|x|^2}{t^{2\beta}}\right)_+^{1/(\gamma-1)}.
\]
\end{example}

\begin{remark}[Finite propagation speed for the porous medium equation]
\label{rem:porous-medium-finite-propagation-speed}
\uses{ex:porous-medium-scaling-invariant-barenblatt}
\dcref{ch4:4.2.2-rem-1}
\group{Similarity Solutions}

The Barenblatt profile has compact support at each $t>0$; degenerate diffusion
therefore gives finite propagation speed for compactly supported nonnegative
data.
\end{remark}

\section{Transform Methods}
\label{sec:transform-methods}

\subsection{Fourier Transform}
\label{sec:fourier-transform}

All functions in this subsection may be complex-valued.

\begin{definition}[Fourier transform on $L^1$]
\label{def:fourier-transform-l1}
\dcref{ch4:4.3.1-def-1}
\group{Fourier Transform}
For $u\in L^1(\R^n)$ define
\[
\hat u(y)=(2\pi)^{-n/2}\int_{\R^n}e^{-ix\cdot y}u(x)\,dx,
\qquad
\check u(y)=(2\pi)^{-n/2}\int_{\R^n}e^{ix\cdot y}u(x)\,dx.
\]
\end{definition}

\begin{theorem}[Plancherel's Theorem]
\label{thm:plancherel-theorem}
\uses{def:fourier-transform-l1, ex:bessel-potentials, lem:integral-fundamental-solution-heat}
\dcref{ch4:4.3.1-thm-1}
\group{Fourier Transform}

For $u\in L^1(\R^n)\cap L^2(\R^n)$,
\[
\hat u,\check u\in L^2(\R^n),\qquad
\|\hat u\|_2=\|\check u\|_2=\|u\|_2.
\]
\end{theorem}

\begin{proof}
\sketch Apply the Gaussian regularization identity to $u*\bar u(-\cdot)$,
use the Fourier convolution formula, and let the Gaussian parameter tend to
zero. This proves $\|\hat u\|_2=\|u\|_2$; the inverse transform is analogous.
\end{proof}

\begin{definition}[Fourier transform on $L^2$]
\label{def:fourier-transform-l2}
\uses{thm:plancherel-theorem}
\dcref{ch4:4.3.1-def-2}
\group{Fourier Transform}

For $u\in L^2(\R^n)$, choose $u_k\in L^1\cap L^2$ with $u_k\to u$ in $L^2$.
Plancherel makes $\hat u_k$ and $\check u_k$ Cauchy in $L^2$; their limits
define $\hat u$ and $\check u$, independently of the approximating sequence.
\end{definition}

\begin{theorem}[Properties of Fourier transform]
\label{thm:fourier-transform-properties}
\uses{def:fourier-transform-l2, thm:plancherel-theorem}
\dcref{ch4:4.3.1-thm-2}
\group{Fourier Transform}

For $u,v\in L^2(\R^n)$,
\[
\int u\bar v\,dx=\int\hat u\,\overline{\hat v}\,dy,
\quad \widehat{D^\alpha u}=(iy)^\alpha\hat u,
\quad (u*v)^\wedge=(2\pi)^{n/2}\hat u\hat v,
\quad u=(\hat u)^\vee,
\]
whenever the indicated derivatives lie in $L^2$.
\end{theorem}

\begin{proof}
\sketch Polarization of Plancherel gives the inner-product formula.
Integration by parts proves the derivative rule for smooth compactly supported
functions and density extends it. Fubini gives convolution; Gaussian
approximate identities give Fourier inversion almost everywhere.
\end{proof}

\begin{example}[Bessel potentials]
\label{ex:bessel-potentials}
\uses{thm:fourier-transform-properties}
\dcref{ch4:4.3.1-ex-1}
\group{Fourier Transform}

For $-\Delta u+u=f$,
\[
\hat u(y)=\frac{\hat f(y)}{1+|y|^2},\qquad
u=\frac{f*B}{(2\pi)^{n/2}},
\]
where
\[
B(x)=\frac1{2\pi^{n/2}}\int_0^\infty t^{-n/2}e^{-t-|x|^2/(4t)}\,dt.
\]
\end{example}

\begin{example}[Fundamental solution of the heat equation, via Fourier transform]
\label{ex:fourier-transform-heat-fundamental-solution}
\uses{thm:fourier-transform-properties, ex:bessel-potentials, thm:heat-equation-initial-value-solution}
\dcref{ch4:4.3.1-ex-2}
\group{Fourier Transform}

Fourier transformation in $x$ reduces the heat initial-value problem to
$\hat u_t+|y|^2\hat u=0$, so
\[
u(x,t)=\frac1{(4\pi t)^{n/2}}\int_{\R^n}e^{-|x-y|^2/(4t)}g(y)\,dy.
\]
\end{example}

\begin{example}[Fundamental solution of Schr\"odinger's equation]
\label{ex:schrodinger-fundamental-solution}
\uses{thm:plancherel-theorem, ex:fourier-transform-heat-fundamental-solution}
\dcref{ch4:4.3.1-ex-3}
\group{Fourier Transform}

For $iu_t+\Delta u=0$, formally continue the heat kernel to $it$:
\[
u(x,t)=\frac1{(4\pi it)^{n/2}}\int_{\R^n}e^{i|x-y|^2/(4t)}g(y)\,dy.
\]
This solves the equation for suitable $g$, and Plancherel gives
$\|u(\cdot,t)\|_2=\|g\|_2$, extending the formula from $L^1\cap L^2$ to
$L^2$.
\end{example}

\begin{remark}[Fundamental solution of Schr\"odinger's equation and time reversibility]
\label{rem:schrodinger-time-reversibility}
\uses{ex:schrodinger-fundamental-solution, thm:backwards-uniqueness-heat-equation}
\dcref{ch4:4.3.1-rem-1}
\group{Fourier Transform}

The fundamental solution
\[
\Psi(x,t)=(4\pi it)^{-n/2}e^{i|x|^2/(4t)},\qquad t\ne0,
\]
defines $u=g*\Psi$ for positive and negative times; Schrodinger evolution is
time reversible, unlike heat evolution.
\end{remark}

\begin{example}[The wave equation via Fourier transform]
\label{ex:fourier-transform-wave-equation}
\uses{thm:fourier-transform-properties, ex:exponential-plane-wave-solutions}
\dcref{ch4:4.3.1-ex-4}
\group{Fourier Transform}

For $u_{tt}-\Delta u=0$, $u(\cdot,0)=g$, $u_t(\cdot,0)=0$,
\[
\hat u(y,t)=\frac{\hat g(y)}2(e^{it|y|}+e^{-it|y|}),
\]
and hence
\[
u(x,t)=\frac1{(2\pi)^{n/2}}\int_{\R^n}\frac{\hat g(y)}2
\left(e^{i(x\cdot y+t|y|)}+e^{i(x\cdot y-t|y|)}\right)dy.
\]
\end{example}

\begin{example}[Telegraph equation]
\label{ex:telegraph-equation-fourier-transform}
\uses{thm:fourier-transform-properties}
\dcref{ch4:4.3.1-ex-5}
\group{Fourier Transform}

For $u_{tt}+2du_t-u_{xx}=0$, Fourier modes satisfy
$\hat u_{tt}+2d\hat u_t+|y|^2\hat u=0$. With
$\gamma=(d^2-|y|^2)^{1/2}$ and $\delta=(|y|^2-d^2)^{1/2}$,
\[
\hat u=e^{-dt}\begin{cases}
\beta_1e^{\gamma t}+\beta_2e^{-\gamma t},&|y|\le d,\\
\beta_1e^{i\delta t}+\beta_2e^{-i\delta t},&|y|\ge d,
\end{cases}
\]
where $\beta_1+\beta_2=\hat g$ and the derivative at $t=0$ equals $\hat h$.
\end{example}

\subsection{Laplace Transform}
\label{sec:laplace-transform}

\begin{definition}[Laplace transform]
\label{def:laplace-transform}
\dcref{ch4:4.3.2-def-1}
\group{Laplace Transform}
For $u\in L^1(\R_+)$, define
\[
u^\#(s)=\int_0^\infty e^{-st}u(t)\,dt,\qquad s\ge0.
\]
\end{definition}

\begin{example}[Resolvents and Laplace transform]
\label{ex:resolvent-laplace-transform}
\uses{def:laplace-transform, ex:bessel-potentials}
\dcref{ch4:4.3.2-ex-1}
\group{Laplace Transform}

For $v_t-\Delta v=0$ with $v(\cdot,0)=f$, the time transform
$u=v^\#(\cdot,s)$ satisfies
\[
-\Delta u+su=f.
\]
\end{example}

\begin{example}[Wave equation from the heat equation, via Laplace transform]
\label{ex:wave-equation-from-heat-laplace-transform}
\uses{def:laplace-transform, thm:wave-equation-odd-dimensions-solution}
\dcref{ch4:4.3.2-ex-2}
\group{Laplace Transform}

For odd $n$, extend the zero-velocity wave solution evenly to negative time and
set
\[
v(x,t)=(4\pi t)^{-1/2}\int_\R e^{-s^2/(4t)}u(x,s)\,ds.
\]
Then $v_t-\Delta v=0$, $v(\cdot,0)=g$. Comparing its heat-kernel formula and
integrating by parts gives, for $n=2k+1$,
\[
u(x,t)=\frac1{\gamma_n}\frac\partial{\partial t}
\left(\frac1t\frac\partial{\partial t}\right)^{(n-3)/2}
\left(t^{n-2}\int_{\partial B(x,t)}g\,dS\right),
\]
where $\gamma_n=(n-2)(n-4)\cdots3$.
The normalization uses the Gamma-function recursion from \cite{Rudin1976}.
\end{example}

\section{Converting Nonlinear into Linear PDE}
\label{sec:converting-nonlinear-into-linear-pde}

\subsection{Hopf--Cole Transformation}
\label{sec:hopf-cole-transformation}

\begin{example}[The Hopf--Cole transformation for a quadratic-nonlinearity parabolic PDE]
\label{ex:hopf-cole-quadratic-nonlinearity}
\uses{ex:fourier-transform-heat-fundamental-solution}
\dcref{ch4:4.4.1-ex-1}
\group{Nonlinear to Linear PDE}

For
\[
u_t-a\Delta u+b|Du|^2=0,\qquad u(x,0)=g(x),\quad a>0,
\]
set $w=\phi(u)$ and choose $a\phi''+b\phi'=0$, namely
$w=e^{-bu/a}$. Then $w_t-a\Delta w=0$ with $w(x,0)=e^{-bg(x)/a}$, so
\[
u(x,t)=-\frac ab\log\left[(4\pi at)^{-n/2}
\int_{\R^n}e^{-|x-y|^2/(4at)-bg(y)/a}\,dy\right].
\]
\end{example}

\begin{example}[Viscous Burgers' equation via the Hopf--Cole transformation]
\label{ex:viscous-burgers-hopf-cole}
\uses{ex:hopf-cole-quadratic-nonlinearity, def:integral-solution-conservation-law}
\dcref{ch4:4.4.1-ex-2}
\group{Nonlinear to Linear PDE}

For viscous Burgers,
\[
u_t-au_{xx}+uu_x=0,\qquad u(x,0)=g(x),
\]
set $w(x,t)=\int_{-\infty}^xu(y,t)\,dy$ and
$h(x)=\int_{-\infty}^xg(y)\,dy$. Then
$w_t-aw_{xx}+\frac12w_x^2=0$, and Hopf--Cole gives
\[
u(x,t)=\frac{\displaystyle\int_\R\frac{x-y}{t}
e^{-|x-y|^2/(4at)-h(y)/(2a)}\,dy}
{\displaystyle\int_\R e^{-|x-y|^2/(4at)-h(y)/(2a)}\,dy}.
\]
\end{example}

\subsection{Potential Functions}
\label{sec:potential-functions}

\begin{example}[Potential functions for Euler's equations]
\label{ex:euler-equations-potential-function}
\dcref{ch4:4.4.2-ex-1}
\group{Nonlinear to Linear PDE}
For incompressible Euler flow
\[
\mathbf u_t+\mathbf u\cdot D\mathbf u=-Dp+\mathbf f,\qquad
\Div\mathbf u=0,
\]
assume $\Div\mathbf g=0$, $\mathbf f=Dh$, and seek $\mathbf u=Dv$ with
$Dv(\cdot,0)=\mathbf g$. Then $\Delta v=0$ and
\[
v_t+\frac12|Dv|^2+p=h,
\]
which recovers the pressure once the harmonic potential $v$ is known.
\end{example}

\subsection{Hodograph and Legendre Transforms}
\label{sec:hodograph-legendre-transforms}

\begin{example}[Hodograph transform for steady irrotational fluid flow]
\label{ex:hodograph-transform-irrotational-flow}
\dcref{ch4:4.4.3-ex-1}
\group{Nonlinear to Linear PDE}
For the steady two-dimensional irrotational flow system, invert
$\mathbf u=\mathbf u(x)$ to $\mathbf x=\mathbf x(\mathbf u)$ wherever
\[
J=\det D_x\mathbf u\ne0.
\]
The derivative identities
\[
u^2_{x_2}=Jx^1_{u_1},\quad u^2_{x_1}=-Jx^2_{u_1},\quad
u^1_{x_2}=-Jx^1_{u_2},\quad u^1_{x_1}=Jx^2_{u_2}
\]
convert the quasilinear system to the linear system
\[
\begin{cases}
(\sigma^2-u_1^2)x^2_{u_2}+u_1u_2(x^2_{u_2}+x^2_{u_1})
+(\sigma^2-u_2^2)x^1_{u_1}=0,\\
x^1_{u_2}-x^2_{u_1}=0.
\end{cases}
\]
\end{example}

\begin{remark}[Simplifying the hodograph system via a potential function]
\label{rem:hodograph-potential-simplification}
\uses{ex:hodograph-transform-irrotational-flow}
\dcref{ch4:4.4.3-rem-1}
\group{Nonlinear to Linear PDE}

Writing $x^1=z_{u_1}$ and $x^2=z_{u_2}$ reduces the hodograph system to
\[
(\sigma^2-u_1^2)z_{u_2u_2}+2u_1u_2z_{u_1u_2}
+(\sigma^2-u_2^2)z_{u_1u_1}=0.
\]
\end{remark}

\begin{example}[Legendre transform for the minimal surface equation]
\label{ex:legendre-transform-minimal-surface}
\uses{def:legendre-transform}
\dcref{ch4:4.4.3-ex-2}
\group{Nonlinear to Linear PDE}

For the two-dimensional minimal-surface equation, invert
$p^i=u_{x_i}$ where $J=\det D^2u\ne0$, and define
\[
v(p)=x(p)\cdot p-u(x(p)).
\]
Then
\[
u_{x_1x_1}=Jv_{p_2p_2},\quad u_{x_1x_2}=-Jv_{p_1p_2},\quad
u_{x_2x_2}=Jv_{p_1p_1},
\]
and the minimal-surface equation becomes the linear equation
\[
(1+p_2^2)v_{p_2p_2}+2p_1p_2v_{p_1p_2}+(1+p_1^2)v_{p_1p_1}=0.
\]
\end{example}

\begin{remark}[Practical difficulty of the hodograph and Legendre techniques]
\label{rem:hodograph-legendre-practical-difficulty}
\uses{ex:hodograph-transform-irrotational-flow, ex:legendre-transform-minimal-surface}
\dcref{ch4:4.4.3-rem-2}
\group{Nonlinear to Linear PDE}

These transforms linearize selected equations but generally make prescribed
boundary conditions difficult to express.
\end{remark}

\section{Asymptotics}
\label{sec:asymptotics}

\subsection{Singular Perturbations}
\label{sec:singular-perturbations}

\begin{definition}[Singular perturbation]
\label{def:singular-perturbation}
\dcref{ch4:4.5.1-def-1}
\group{Asymptotics}
A singular perturbation adds $\varepsilon$ times a higher-order term to a PDE;
the limit $\varepsilon\to0$ can change the analytic character of solutions.
\end{definition}

\begin{example}[Transport and small diffusion]
\label{ex:transport-small-diffusion-singular-perturbation}
\uses{def:singular-perturbation}
\dcref{ch4:4.5.1-ex-1}
\group{Asymptotics}

Let a smooth planar velocity field $\mathbf b$ carry dye injected at the
origin. The unperturbed density solves
\[
\Div(u\mathbf b)=\delta_0.
\]
Along $\dot{\mathbf x}=\mathbf b(\mathbf x)$, with speed
$\sigma(t)=|\mathbf b(\mathbf x(t))|$, introduce normal coordinates
$x=\mathbf x(t)+y\nu$; the Jacobian is $\sigma(1-\kappa y)$. The weak density is
\[
u=\frac{\delta(y)}{\sigma(t)}.
\]
For the perturbation $-\varepsilon\Delta u^\varepsilon+\Div(u^\varepsilon\mathbf b)=\delta_0$,
set $z=\varepsilon^{-1/2}y$, $v^\varepsilon=\varepsilon^{1/2}u^\varepsilon$.
If
$\mathbf b=\sigma\tau+(\alpha\tau+\beta\nu)y+O(y^2)$ and
$v^\varepsilon\to v$, the formal limit is
\[
v_t-v_{zz}+(\beta zv)_z+\frac{\dot\sigma}{\sigma}v=0,
\qquad
v(z,0)=\frac{\beta(z)}{\sigma(0)},
\]
with $u^\varepsilon=\varepsilon^{-1/2}(v+o(1))$.
\end{example}

\subsection{Laplace's Method}
\label{sec:laplaces-method}

\begin{example}[Vanishing viscosity method for Burgers' equation]
\label{ex:vanishing-viscosity-burgers}
\uses{ex:viscous-burgers-hopf-cole, thm:lax-oleinik-formula}
\dcref{ch4:4.5.2-ex-1}
\group{Asymptotics}

For viscous Burgers with viscosity $\varepsilon$,
\[
u_t^\varepsilon+u^\varepsilon u_x^\varepsilon-\varepsilon u_{xx}^\varepsilon=0,
\]
the Hopf--Cole formula is
\[
u^\varepsilon(x,t)=
\frac{\int_\R \frac{x-y}{t}e^{-K(x,y,t)/\varepsilon}\,dy}
{\int_\R e^{-K(x,y,t)/\varepsilon}\,dy},
\qquad
K(x,y,t)=\frac{|x-y|^2}{2t}+h(y).
\]
At every $x$ where $K(\cdot,x,t)$ has a unique minimizer $y(x,t)$,
\[
\lim_{\varepsilon\to0}u^\varepsilon(x,t)=\frac{x-y(x,t)}t,
\]
the Lax--Oleinik entropy solution.
\end{example}

\begin{lemma}[Laplace-method asymptotics for a ratio of Gaussian-type integrals]
\label{lem:laplace-method-asymptotics-lemma}
\dcref{ch4:4.5.2-lem-1}
\group{Asymptotics}
If continuous $k,l:\R\to\R$ have $k$ growing at least quadratically,
$l$ at most linearly, and $k$ has a unique minimizer $y_0$, then
\[
\lim_{\varepsilon\to0}
\frac{\int_\R l(y)e^{-k(y)/\varepsilon}\,dy}
{\int_\R e^{-k(y)/\varepsilon}\,dy}=l(y_0).
\]
\end{lemma}

\begin{proof}
\sketch Normalize $e^{-(k-k(y_0))/\varepsilon}$ to a probability density.
Quadratic growth gives domination and concentration at the unique minimizer;
the weighted averages of the at-most-linear function $l$ therefore converge to
$l(y_0)$.
\end{proof}

\subsection{Geometric Optics, Stationary Phase}
\label{sec:geometric-optics-stationary-phase}

\begin{example}[Oscillating solutions and the geometric optics ansatz]
\label{ex:geometric-optics-oscillating-solutions}
\uses{thm:characteristic-ode-structure}
\dcref{ch4:4.5.3-ex-1}
\group{Asymptotics}

For $u_{tt}-\Delta u=0$, insert
\[
u^\varepsilon=e^{ip^\varepsilon/\varepsilon}a^\varepsilon.
\]
The real part gives
\[
a^\varepsilon\big((p_t^\varepsilon)^2-|Dp^\varepsilon|^2\big)
=\varepsilon^2(a_{tt}^\varepsilon-\Delta a^\varepsilon).
\]
If $p^\varepsilon\to p$ and $a^\varepsilon\to a\ne0$, the limiting phase
satisfies $p_t\pm|Dp|=0$. For a hyperbolic operator with coefficients
$a^{kl}(x)$, the corresponding eikonal equation is
\[
p_t\pm\left(\sum_{k,l}a^{kl}p_{x_k}p_{x_l}\right)^{1/2}=0.
\]
\end{example}

\begin{example}[Stationary phase for the wave equation]
\label{ex:stationary-phase-wave-equation}
\uses{ex:fourier-transform-wave-equation}
\dcref{ch4:4.5.3-ex-2}
\group{Asymptotics}

For rapidly oscillatory initial data $g^\varepsilon(x)=a(x)e^{ip(x)/\varepsilon}$,
the Fourier wave formula becomes
\[
u^\varepsilon=\frac12(I_+^\varepsilon+I_-^\varepsilon),\qquad
I_\pm^\varepsilon=(2\pi\varepsilon)^{-n}
\iint a(z)e^{i\phi_\pm(x,y,z,t)/\varepsilon}\,dy\,dz,
\]
where
\[
\phi_\pm=(x-z)\cdot y\pm t|y|+p(z).
\]
\end{example}

\begin{lemma}[Asymptotics for linear phase]
\label{lem:stationary-phase-linear-asymptotics}
\dcref{ch4:4.5.3-lem-1}
\group{Asymptotics}
If $a\in C_c^\infty(\R^n)$ and $p\ne0$, then for every integer $m\ge1$,
\[
\int_{\R^n}e^{ip\cdot y/\varepsilon}a(y)\,dy=O(\varepsilon^m).
\]
\end{lemma}

\begin{proof}
\sketch Choose a coordinate with $p_1\ne0$ and integrate by parts $m$ times
using
\[
e^{ip\cdot y/\varepsilon}=(\varepsilon/(ip_1))\partial_{y_1}
e^{ip\cdot y/\varepsilon}.
\]
\end{proof}

\begin{lemma}[Asymptotics for quadratic phase]
\label{lem:stationary-phase-quadratic-asymptotics}
\dcref{ch4:4.5.3-lem-2}
\group{Asymptotics}
If $a\in C_c^\infty(\R^n)$ and $A$ is real, symmetric, and nonsingular, then
\[
\frac1{(2\pi)^{n/2}}\int e^{iy\cdot Ay/(2\varepsilon)}a(y)\,dy
=\frac{e^{i\pi\operatorname{sgn}(A)/4}}{|\det A|^{1/2}}
\big(a(0)+O(\varepsilon)\big).
\]
\end{lemma}

\begin{proof}
\sketch Add Gaussian damping, diagonalize $A$, and evaluate the resulting
one-dimensional complex Gaussian integrals by contour deformation. Fourier
inversion of the resulting oscillatory Gaussian and Taylor expansion of $a$
give the $O(\varepsilon)$ remainder.
\end{proof}

\begin{lemma}[Changing coordinates near a critical point]
\label{lem:stationary-phase-changing-coordinates}
\dcref{ch4:4.5.3-lem-3}
\group{Asymptotics}
If $\phi$ is smooth, then:
\begin{enumerate}
\item If $D\phi(0)\ne0$, there is a local diffeomorphism $\Phi$ with
$\Phi(0)=0$, $D\Phi(0)=I$, and
$\phi(\Phi(x))=\phi(0)+D\phi(0)\cdot x$.
\item If $D\phi(0)=0$ and $\det D^2\phi(0)\ne0$, there is such a $\Phi$ with
$\phi(\Phi(x))=\phi(0)+\frac12x\cdot D^2\phi(0)x$.
\end{enumerate}
\end{lemma}

\begin{proof}
\sketch The first statement follows from the implicit function theorem after
using $D\phi(0)$ as one coordinate. For the second, Taylor-expand
$\phi(x)=\phi(0)+\frac12x\cdot A(x)x$, diagonalize $A(0)$, and successively
complete squares; each coordinate change has derivative $I$ at the origin.
\end{proof}

\begin{proposition}[The stationary phase asymptotic formula]
\label{prop:stationary-phase-asymptotic-formula}
\uses{lem:stationary-phase-linear-asymptotics, lem:stationary-phase-quadratic-asymptotics, lem:stationary-phase-changing-coordinates}
\dcref{ch4:4.5.3-prop-1}
\group{Asymptotics}

For $I_\varepsilon=\int e^{i\phi(y)/\varepsilon}a(y)\,dy$, assume the
critical points $y_1,\ldots,y_N$ in $\operatorname{spt}a$ are the only zeros
of $D\phi$ and $D^2\phi(y_k)$ is nonsingular. Then
\[
I_\varepsilon=(2\pi\varepsilon)^{n/2}\sum_{k=1}^N
\frac{e^{i\phi(y_k)/\varepsilon}e^{i\pi\operatorname{sgn}(D^2\phi(y_k))/4}}
{|\det D^2\phi(y_k)|^{1/2}}\big(a(y_k)+O(\varepsilon)\big).
\]
\end{proposition}

\begin{proof}
\sketch A partition of unity separates small neighborhoods of the critical
points from a region where $D\phi\ne0$. The latter is $O(\varepsilon^m)$ by
the linear-phase lemma; the former are reduced to quadratic phase by the
coordinate lemma and evaluated by the quadratic-phase lemma.
\end{proof}

\begin{example}[Stationary phase for the wave equation, continued]
\label{ex:stationary-phase-wave-equation-continued}
\uses{ex:stationary-phase-wave-equation, prop:stationary-phase-asymptotic-formula}
\dcref{ch4:4.5.3-ex-3}
\group{Asymptotics}

For $\phi_+=(x-z)\cdot y+t|y|+p(z)$, the stationary set is
\[
S^+=\left\{(y,z):x=z-t\frac{Dp(z)}{|Dp(z)|},\ y=Dp(z)\right\}.
\]
At a stationary point,
\[
\det D^2_{y,z}\phi_+=(-1)^n\prod_{i=1}^{n-1}(1-t\kappa_i(z)),
\]
where $\kappa_i$ are principal curvatures of the level set of $p$. For small
$t$, the implicit function theorem selects a unique stationary point and the
stationary-phase proposition gives the corresponding asymptotic for
$I_+^\varepsilon$; $I_-^\varepsilon$ is analogous.
\end{example}

\begin{remark}[Optics and stationary phase]
\label{rem:optics-stationary-phase-connection}
\uses{ex:geometric-optics-oscillating-solutions, ex:stationary-phase-wave-equation-continued, thm:characteristic-ode-structure}
\dcref{ch4:4.5.3-rem-1}
\group{Asymptotics}

The eikonal equations $p_t\pm|Dp|=0$ have characteristics
\[
\dot x=\mp p/|p|,\qquad \dot p=0.
\]
For small time, these straight rays are exactly the stationary sets selected
by the phases $\phi_\pm$.
\end{remark}

\subsection{Homogenization}
\label{sec:homogenization}

\begin{example}[Periodic homogenization of an elliptic equation]
\label{ex:periodic-homogenization-elliptic}
\dcref{ch4:4.5.4-ex-1}
\group{Asymptotics}
For uniformly elliptic $Q$-periodic coefficients $a^{ij}(y)$, consider
\[
-\sum_{i,j}(a^{ij}(x/\varepsilon)u^\varepsilon_{x_i})_{x_j}=f\text{ in }U,
\qquad u^\varepsilon|_{\partial U}=0.
\]
The formal expansion
\[
u^\varepsilon=u_0(x,x/\varepsilon)+\varepsilon u_1(x,x/\varepsilon)+\cdots
\]
and $L=\varepsilon^{-2}L_1+\varepsilon^{-1}L_2+L_3$ imply that $u_0=u(x)$
is independent of the fast variable. Correctors $\chi^i$, periodic solutions
of
\[
L_1\chi^i=-\sum_j a^{ij}(y)u_j,
\]
give $u_1=-\sum_i\chi^i(y)u_{x_i}+\tilde u_1(x)$. Solvability of the next
cell problem yields the homogenized equation
\[
-\sum_{i,j}\bar a^{ij}u_{x_ix_j}=f\text{ in }U,\qquad u|_{\partial U}=0,
\]
with
\[
\bar a^{ij}=\int_Q\left(a^{ij}(y)-\sum_k a^{jk}(y)\chi^i_{y_k}(y)\right)dy.
\]
\end{example}

\section{Power Series}
\label{sec:power-series}

\subsection{Noncharacteristic Surfaces}
\label{sec:noncharacteristic-surfaces}

Consider the $k$th-order quasilinear equation
\[
\sum_{|\alpha|=k}a_\alpha(D^{k-1}u,\ldots,u,x)D^\alpha u+a_0(D^{k-1}u,\ldots,u,x)=0
\]
near a smooth hypersurface $\Gamma$ with unit normal $\nu$.

\begin{definition}[Normal derivative]
\label{def:normal-derivative-notation}
\dcref{ch4:4.6.1-def-1}
\group{Power Series Methods}
Define
\[
\frac{\partial^ju}{\partial\nu^j}:=
\sum_{|\alpha|=j}D^\alpha u\,\nu^\alpha.
\]
\end{definition}

\begin{definition}[Cauchy problem and Cauchy data]
\label{def:cauchy-problem-cauchy-data}
\uses{def:normal-derivative-notation}
\dcref{ch4:4.6.1-def-2}
\group{Power Series Methods}

Given $g_0,\ldots,g_{k-1}:\Gamma\to\R$, the Cauchy problem prescribes
\[
\frac{\partial^ju}{\partial\nu^j}=g_j\quad\text{on }\Gamma,
\qquad j=0,\ldots,k-1,
\]
along with the PDE.
\end{definition}

\begin{definition}[Noncharacteristic plane]
\label{def:noncharacteristic-plane-flat}
\dcref{ch4:4.6.1-def-3}
\group{Power Series Methods}
For the plane $\Gamma=\{x_n=0\}$, call the plane noncharacteristic if
\[
a_{(0,\ldots,0,k)}\ne0
\]
for all coefficient arguments.
\end{definition}

If the plane is noncharacteristic, tangential differentiation of the Cauchy
data determines all derivatives through order $k-1$ and all order-$k$
derivatives except $u_{x_n^k}$. The PDE determines the missing term by
\[
u_{x_n^k}=-\frac1{a_{(0,\ldots,0,k)}}
\left[\sum_{|\alpha|=k,\,\alpha\ne(0,\ldots,0,k)}a_\alpha D^\alpha u+a_0\right].
\]
Differentiating the PDE repeatedly then determines every derivative on the
plane.

\begin{definition}[Noncharacteristic surface]
\label{def:noncharacteristic-surface}
\uses{def:cauchy-problem-cauchy-data}
\dcref{ch4:4.6.1-def-4}
\group{Power Series Methods}

An arbitrary smooth hypersurface $\Gamma$ is noncharacteristic if
\[
\sum_{|\alpha|=k}a_\alpha\nu^\alpha\ne0\quad\text{on }\Gamma
\]
for all coefficient arguments.
\end{definition}

\begin{theorem}[Cauchy data and noncharacteristic surfaces]
\label{thm:cauchy-data-noncharacteristic-surfaces}
\uses{def:noncharacteristic-surface, def:cauchy-problem-cauchy-data, def:noncharacteristic-plane-flat}
\dcref{ch4:4.6.1-thm-1}
\group{Power Series Methods}

If $\Gamma$ is noncharacteristic, every derivative of a smooth solution along
$\Gamma$ is uniquely determined by the Cauchy data and the PDE coefficients.
\end{theorem}

\begin{proof}
\sketch Flatten $\Gamma$ by a local diffeomorphism. The transformed highest
normal coefficient is a nonzero multiple of
$\sum_{|\alpha|=k}a_\alpha\nu^\alpha$, so the flat-plane induction applies to
the transformed Cauchy data. Transform the resulting derivatives back.
\end{proof}

\begin{remark}[Noncharacteristic condition via a level set function]
\label{rem:noncharacteristic-condition-level-set}
\uses{def:noncharacteristic-surface}
\dcref{ch4:4.6.1-rem-1}
\group{Power Series Methods}

If $\Gamma=\{w=0\}$ and $Dw\ne0$ on $\Gamma$, then the condition is
\[
\sum_{|\alpha|=k}a_\alpha(Dw)^\alpha\ne0\quad\text{on }\Gamma.
\]
\end{remark}

\subsection{Real Analytic Functions}
\label{sec:real-analytic-functions}

\begin{definition}[Real analytic function]
\label{def:real-analytic-function}
\dcref{ch4:4.6.2-def-1}
\group{Power Series Methods}
A function $f:\R^n\to\R$ is real analytic near $x_0$ if
\[
f(x)=\sum_\alpha f_\alpha(x-x_0)^\alpha
\]
for $|x-x_0|<r$ and some $r>0$.
\end{definition}

\begin{remark}[Basic properties of real analytic functions]
\label{rem:real-analytic-function-remarks}
\uses{def:real-analytic-function}
\dcref{ch4:4.6.2-rem-1}
\group{Power Series Methods}

Analytic functions are $C^\infty$ and
\[
f_\alpha=\frac{D^\alpha f(x_0)}{\alpha!},\qquad
f(x)=\sum_\alpha\frac{D^\alpha f(x_0)}{\alpha!}(x-x_0)^\alpha.
\]
\end{remark}

\begin{example}[Power series for a simple rational function]
\label{ex:real-analytic-example-rational-function}
\uses{def:real-analytic-function}
\dcref{ch4:4.6.2-ex-1}
\group{Power Series Methods}

For $|x|<r/\sqrt n$,
\[
\frac r{r-(x_1+\cdots+x_n)}
=\sum_\alpha\frac{|\alpha|!}{r^{|\alpha|}\alpha!}x^\alpha,
\]
absolutely, by the multinomial theorem and
$|x_1|+\cdots+|x_n|\le\sqrt n|x|$.
\end{example}

\begin{definition}[Majorizes]
\label{def:power-series-majorizes}
\dcref{ch4:4.6.2-def-2}
\group{Power Series Methods}
For $f=\sum f_\alpha x^\alpha$ and $g=\sum g_\alpha x^\alpha$, write
$g\gg f$ if $g_\alpha\ge|f_\alpha|$ for every multiindex $\alpha$.
\end{definition}

\begin{lemma}[Majorants]
\label{lem:power-series-majorants}
\uses{def:power-series-majorizes, ex:real-analytic-example-rational-function}
\dcref{ch4:4.6.2-lem-1}
\group{Power Series Methods}

If $g\gg f$ and $g$ converges for $|x|<r$, then $f$ converges there. If $f$
converges for $|x|<r$ and $0<s\sqrt n<r$, then $f$ has a majorant on
$|x|<s\sqrt n$.
\end{lemma}

\begin{proof}
\sketch The coefficient inequality gives absolute convergence of $f$ wherever
$g$ converges. For the converse, evaluate the convergent series at
$y=s(1,\ldots,1)$ to bound $|f_\alpha|$, then majorize by
$Cs/[s-(x_1+\cdots+x_n)]$.
\end{proof}

\begin{remark}[Vector-valued power series notation]
\label{rem:vector-valued-power-series-notation}
\uses{def:power-series-majorizes}
\dcref{ch4:4.6.2-rem-2}
\group{Power Series Methods}

For vectors of power series, write $\mathbf g\gg\mathbf f$ when
$g^k\gg f^k$ for every component.
\end{remark}

\subsection{Cauchy--Kovalevskaya Theorem}
\label{sec:cauchy-kovalevskaya-theorem}

After flattening an analytic noncharacteristic surface and subtracting analytic
Cauchy data, the $k$th-order problem reduces to
\[
\mathbf u_{x_n}=\sum_{j=1}^{n-1}\mathbf B_j(\mathbf u,x')\mathbf u_{x_j}
+\mathbf c(\mathbf u,x'),\qquad \mathbf u|_{x_n=0}=0,
\]
where the coefficients are analytic and independent of $x_n$.

\begin{theorem}[Cauchy--Kovalevskaya Theorem]
\label{thm:cauchy-kovalevskaya}
\uses{def:cauchy-problem-cauchy-data, thm:cauchy-data-noncharacteristic-surfaces, lem:power-series-majorants, rem:vector-valued-power-series-notation, def:real-analytic-function}
\dcref{ch4:4.6.3-thm-1}
\group{Power Series Methods}

If $\mathbf B_j$ and $\mathbf c$ are real analytic, there exist $r>0$ and a
real analytic solution
\[
\mathbf u(x)=\sum_\alpha\mathbf u_\alpha x^\alpha
\]
of the reduced system near the origin.
\end{theorem}

\begin{proof}
\sketch Differentiate the system recursively to express each Taylor
coefficient as a polynomial with nonnegative coefficients in coefficient
derivatives and lower normal-order coefficients. Choose analytic majorants
for $\mathbf B_j$ and $\mathbf c$; induction gives a componentwise majorant
for the unknown series. A scalar explicit majorant solves the resulting model
system and is analytic in a small ball, proving convergence. The convergent
series satisfies the system because its Taylor expansions agree.
\end{proof}

\begin{remark}[The Cauchy--Kovalevskaya Theorem for fully nonlinear equations]
\label{rem:cauchy-kovalevskaya-fully-nonlinear}
\uses{thm:cauchy-kovalevskaya}
\dcref{ch4:4.6.3-rem-1}
\group{Power Series Methods}

The theorem extends to fully nonlinear analytic PDE; see \cite{Folland1976}.
\end{remark}

\section{References}
\label{sec:ch4-references}

\noindent\cite{Pinsky1991,Pinsky1991,Strauss1992,ThoeZachmanoglou1986,Weinberger1965,Olver1986,SteinWeiss1971,Stein1993,Hormander1983,Rauch1992,Treves1975,PinskyTaylor1997,CourantHilbert1962,Hormander1971,BensoussanLionsPapanicolaou1978,Folland1976,John1982,DiBenedetto1995,Mikhailov1978}
